\section{相対smoothnessと三次周期性の判定}
\label{sec:relative-saturation-reduction}
本節では滑らかな可換cubicを含むPoisson量子化族の相対saturationを、
全超曲面圏の対角の有限証明書へ還元する。
相対Beilinson仮定(RB1)--(RB2)、相対例外列、相対Orlov比較は
元論文の量子化定理からは自動的に従わない。
gluingに関する以下の命題は既知の一般原理の相対版であり、
新規な分類定理とは主張しない\cite[Theorem 3.25]{OrlovGlue}。

\subsection{例外列を除く前後の相対smoothness}
\begin{proposition}[相対smoothnessの同値]
\label{prop:saturation-reduction}
$R$ を可換Noether環、$\mathcal C$ を小さいidempotent completeな
$R$-linear dg圏とし、
$\RHom_{\mathcal C}(X,Y)\in\Perf(R)$ が任意の対象$X,Y$に対し成り立つとする。
有限個の相対例外対象 $E_1,\ldots,E_r$ と
dg半直交分解
\[
\mathcal C=\langle\mathcal K,E_1,\ldots,E_r\rangle
\]
が与えられ、さらに$\mathcal K$がcompact generatorを持つならば、
\[
\mathcal C\text{ は }R\text{-smooth}
\quad\Longleftrightarrow\quad
\mathcal K\text{ は }R\text{-smooth}.
\]
ゆえに相対propernessが既に分かっていれば、
両者の相対saturationは同値な条件である。
\end{proposition}
\begin{proof}
例外対象を一つずつ追加する帰納法により一個の場合に帰着する。
$\mathcal K$ のcompact generatorを$G$、例外対象を$E$とする。
この半直交分解のdg enhancementは$\mathcal K$と$\Perf(R)$を
\[
M=\RHom_{\mathcal C}(G,E)
\]
により三角dg gluingしたものとして表される（右左の規約によるoppositeを適宜取る）。
$\Gamma=\REnd_{\mathcal K}(G)$ と置けば、propernessより
$M$は$R$-perfectである。
$\mathcal K$がsmoothならば
$\Gamma\in\Perf(\Gamma^{\mathrm{op}}\otL_R\Gamma)$。
恒等式 $M\simeq M\otL_\Gamma\Gamma$ に対角の有限cell表示を適用すると、
その像は $M\otL_R\Gamma$ の有限回のshift、cone、直和因子から構成される。
$M$が$R$-perfectだから$M$は$\Gamma$-perfect、従って
gluing bimoduleがperfectとなる。三角dg圏の対角の標準triangleにより
$\mathcal C$はsmoothとなる。
逆に$\mathcal C$がsmoothなら、同じgluingの対角triangleを
両成分へ制限して$\mathcal K$の対角がperfectである。
以上は有限cell・導来テンソル・直和因子のみを使うので
$R$上でも成立する。体上の定理は\cite[Theorem 3.25]{OrlovGlue}。
\end{proof}

\begin{corollary}[必要な対角証明書は一つでよい]
\label{cor:only-ambient-diagonal}
量子化 $B_R=A_R/(\Omega_R)$ が研究ノートの相対Beilinson仮定
(RB1)--(RB2)と相対例外列を満たし、
$E=\bigoplus_{i=0}^{N-1}\mathcal O_{B_R}(i)$ がcompact generatorであるとする。
$\Lambda_R=\REnd(E)$は$R$-perfectであり、
その残余圏$\mathcal K_R$は、次の条件と同値に$R$-smoothである：
\[
\Lambda_R\in
\Perf(\Lambda_R^{\mathrm{op}}\otL_R\Lambda_R).
\tag{S}
\]
残余圏のEnd代数$\Gamma_R$と$\Lambda_R$の二つの対角を
別々に持ち上げる必要はない。
\end{corollary}
\begin{proof}
研究ノートの相対proper模型によって$\Perf(\Lambda_R)$はproperで、
$\mathcal K_R$と有限相対例外列による半直交分解を持つ。
命題\ref{prop:saturation-reduction}を適用する。
\end{proof}

\begin{remark}[閉ファイバーのsmoothnessだけでは足りない]
$R=\CC[\hbar]_{(\hbar)}$とし、
$\mathcal C_0\simeq\Perf(X)$が滑らかな可換cubicの圏でも、
その対角の有限cell分解を$R$へ持ち上げることは別の問題である。
$P\in\Perf(\Lambda_R^{\mathrm{op}}\otL_R\Lambda_R)$と
$P\to\Lambda_R$を構成し、$\hbar=0$で擬同型であることを証明できれば、
coneの$R$-perfect性とsupportの閉性により
零を含むZariski開近傍で(S)が成り立つ。
この有限証明書はまだ一般量子化族について構成されていない。
\end{remark}

\subsection{graded matrix factorizationの三次周期性}
\begin{lemma}[中心三次元の形式的周期性]
\label{lem:mf-three-period}
$A$を可換環$R$上の整数次数付き結合代数、
$\Omega\in Z(A)_3$を中心元とする。
有限射影次数付き加群のgraded matrix factorizationの
内部次数シフトを$(1)$とすると、自然なdg自己同値の同型
\[
(1)^3\simeq[2]
\]
があり、任意の基底変更と両立する。
\end{lemma}
\begin{proof}
因子分解を
\[
P_0\xrightarrow{d_0}P_1\xrightarrow{d_1}P_0(3),
\quad
d_1d_0=\Omega,\quad d_0(3)d_1=\Omega
\]
と書く。懸垂$[1]$は二項を入れ替え、片側を$(3)$移し、
二つの微分に負号を付ける。二回の懸垂では負号が消え、
両項および微分が一斉に$(3)$移る。
従って$[2]\simeq(3)=(1)^3$である。
次数移動と符号のみを用いるので基底変更とも可換である。
\end{proof}

\begin{corollary}[相対Orlov比較の下でのrotation]
\label{cor:relative-rotation-period}
$\mathcal K_R$と$\MF^{\mathrm{gr}}(A_R,\Omega_R)$の間に、
基底変更と両立する相対Orlov同値があり、
$(1)$が標準rotation$\rho_R$に対応すると仮定する。
すると$\rho_R$は自己同値で、相対的に
\[
\rho_R^3\simeq[2]
\]
が成立する。この段階では相対smoothnessを別途仮定しない。
\end{corollary}
\begin{proof}
補題\ref{lem:mf-three-period}を相対同値に沿って移す。
\end{proof}

\subsection{相対的な可換復元に必要な追加条件}
相対形式復元を代数的な族の全ファイバーへ拡張する際には、次の四点を区別する：
(RB1)--(RB2)を指定量子化族について証明すること、
全超曲面圏$\Perf(\Lambda_R)$の対角の有限証明書を可換点から持ち上げること、
相対Orlov比較でrotationを同定すること、
そして得られたsmooth proper pairが可換cubicのpairの代数的変形像に入ること。
前の三点が解けても最後の可換モデルの存在は別問題であり、
Zhang twistの障害を無視することはできない。
